\chapter{Architettura}

\section{Panoramica dell'Architettura di Sistema}
L'architettura del Migration Tool è stata definita per essere modulare e scalabile.
Alla base del sistema vi è una separazione netta tra il modulo di input (Importer), il motore per l'elaborazione dei metadati,
e il Generatore di Codice.
Questo approccio garantisce una separazione tra il layout dei dati preesistenti e la struttura del database Infor CloudSuite.

\section{Progettazione del Software}
\subsection{Design dell'Importer e del Parser dei Dati}
L'Importer funge da punto d'ingresso per i dati ed è progettato per interpretare il formato proprietario di Infor Exchange,
analizzando file di esportazione testuali o strutturati.
Estraendo i metadati dai record, l'Importer popola i dizionari \textit{Datafiles} (\texttt{txamg110}) e i
relativi \textit{Fields} (\texttt{txamg111}).

Questo design del caricamento rende possibile supportare sia architetture multitenant cloud, sia installazioni legacy on-premise.

\subsection{Motore di Generazione del Codice}
Il nucleo computazionale è rappresentato dal motore di generazione.
Basato su librerie compilate (DLL come \texttt{txamgdll0012} e \texttt{txamgdll0013}), il modulo interpreta il
dizionario di mapping (definito in \texttt{txamg203} o \textit{MappedFields}) e traduce i legami e le regole di
trasformazione in script nativi Baan C.

Questo meccanismo previene di dover scrivere manualmente codice, garantendo il rispetto dei framework CloudSuite:
il codice viene prima generato, poi compilato via \texttt{bic\_compiler}, ed infine agganciato ai \textit{Batch Runs}.

\subsection{Modulo di Integrazione Ad-hoc e Callbacks}
Mentre buona parte dei flussi vengono standardizzati automaticamente, alcune regole di business richiedono sviluppi
specifici a causa della loro complessità.
Vengono quindi esposte a livello architetturale delle interfacce \textit{callback}.

In appositi \textit{module script}, gli sviluppatori possono iniettare blocchi di codice chiamati asincronamente dall'engine.
Questi script personalizzati usufruiscono degli stessi vantaggi transazionali dell'utility principale.

\section{Scelte Tecnologiche}
\subsection{Linguaggi e Framework}
L'intero ecosistema è sviluppato sfruttando l'Infor Tools Framework, programmando in \textbf{Baan C (sia 3GL che 4GL)}.
In parte sono state adottate le \textit{DAL (Data Access Layer)},
per il restante, sono state utilizzati approcci più classici e a basso livello.

\subsection{Librerie di Supporto e Strumenti di Build}
Tutto il coordinamento della configurazione, logging (\texttt{txamgdll0011}) e le \textit{string utility} si basano sulle librerie
del sistema ERP (DLL native Baan e tool come il \texttt{ttdll string}).

Le build non prevedono toolchain esterni complessi
(come Maven o MSBuild), ma si appoggiano direttamente al toolset interno dell'ERP, il quale gestisce i
legami VRC (Version Release Customer) multitenant per garantirne l'incapsulamento del codice generato al corretto Tenant.

